% !TeX root = ../深度学习笔记.tex
% 由主文档通过 \input 载入；请编译主文档。
\setcounter{chapter}{17} % 保留原书章号。
\chapter{多模态学习与跨模态对齐}
\label{chap:dl-multimodal}
同一对象可以由图像、文字或声音描述，各模态保留的信息并不相同。多模态学习的关键是确定哪些观测相互对应、哪些信息需要融合，以及评价任务是否真正要求模型使用多个来源。

\section{模态、配对数据与任务定义}
\emph{模态（modality）}指数据的表现或采集形式。图文检索以一种模态查询另一种，图像描述以图像生成文字，视觉问答同时读取图像和问题；它们的监督单位和输出空间不同。一个图文配对只说明在给定标注协议下二者有关，文字未必描述了全部图像细节，多张图像也可能对应同一描述。

音频可表示为按采样率记录的波形，短时傅里叶变换把局部时间窗映为频率表示，谱图保留随时间变化的频率信息。采样率、窗口和幅度变换改变可见信息与分辨率，不能仅把预处理当作无关格式转换。视频和声音还涉及时间对齐，整段共同出现不保证每一时刻语义对应。

\section{双编码器与共享表示空间}
双编码器分别计算单位化图像向量 $\mathbf u_i$ 和文本向量 $\mathbf v_j$，二者维度相同。令相似度 $S_{ij}=\mathbf u_i^\mathsf T\mathbf v_j/\tau$，$\tau>0$，对 $B$ 对训练样本使用双向目标
\begin{equation}
\mathcal L=-\frac1{2B}\sum_{i=1}^B\left[\log\frac{e^{S_{ii}}}{\sum_je^{S_{ij}}}+\log\frac{e^{S_{ii}}}{\sum_je^{S_{ji}}}\right].
\label{eq:dl-multimodal-contrast}
\end{equation}
第一项给定图像选文本，第二项给定文本选图像，都是式\eqref{eq:dl-contrastive}的应用。相同维度只保证可计算内积，训练目标才促使匹配语义靠近。

如果一个查询存在多个正确候选，应按实际标注构造多正例目标或允许集合，不能把所有非对角项都强制当负例。检索时可以预先缓存候选表示，查询只需编码一次并排序；这提高候选处理效率，但跨模态细粒度交互受到独立编码形式的限制。

\section{跨模态交互与条件生成}
早期融合先组合较底层表示，后期融合组合各模态的预测或高层特征；交叉注意力则允许一方按当前查询读取另一方内容。例如视觉编码器产生 $N$ 个图像特征，经投影后作为键值，文本位置作为查询，输出再用于预测下一个词元。投影后的通道宽度必须与对应注意力头相容。

另一种方式把图像表示作为前缀，与文本嵌入一起送入序列模型。此时掩码应明确图像位置与文本位置的可见关系，生成损失通常仅作用在目标文本上。模型可能利用语言常识直接猜答案而忽略图像；能接受两种输入并不证明它实际利用了两者，需设置输入依赖的诊断。

\section{跨模态评价与捷径诊断}
对给定候选集合，检索 Recall@$k$ 可定义为正确候选集合与前 $k$ 个结果相交的查询比例；候选规模、重复项和多正例规则改变指标含义。生成任务除了语言质量，还需检查具体对象、属性与关系是否由输入支持，单一文本重叠指标不能覆盖全部视觉一致性。

将图文配对打乱、替换关键属性、遮挡相关区域或移除一种模态，可检验任务表现是否依赖相应信息。这些干预也可能使输入离开原分布，因此应尽量构造自然、受控的对照，并比较单模态基线。原始数据中的重复对象和模板文本可能产生捷径，应按来源或对象划分以匹配目标使用情形。

\section{例题与习题}
\begin{example}[对称图文匹配]
两对样本的相似度分数矩阵为 $\mathbf S=\left(\begin{smallmatrix}2&0\\0&2\end{smallmatrix}\right)$。求式\eqref{eq:dl-multimodal-contrast}的损失。
\end{example}
\begin{solve}
每行和每列的正确配对概率都为 $e^2/(e^2+1)$，四个负对数项相同，平均损失为 $\log(1+e^{-2})\approx0.127$。若所有分数相等，损失为 $\log2$。该计算评价配对判别，不能单独证明模型理解了对象间的空间关系。
\end{solve}
\begin{enumerate}
\item 若两条文本都是同一图像的正确描述，应怎样修改匹配目标和检索评价？
\item 为“物体在桌子的左侧还是右侧”设计一组只改变空间关系的图文对照，说明单模态猜测为何不足。
\item 冻结现成编码器，仅训练投影层完成小规模检索实验，同时报告单模态、打乱配对和正常配对结果。
\end{enumerate}

